\section{Points of degree 5 on \texorpdfstring{$X_0(N)$}{X0(N)}: proof of Theorem A}\label{sec:quintic-x0}

Let $C=X_0(N)$ for one of the thirty levels of Theorem~A, of genus $g$, and $J=\Jac C$. The rational cusp
$\infty$ identifies $\Pic^5(C)$ with $J$. For $r\ge0$ let $W^r_5\subseteq\Pic^5(C)$ be the locus of classes
with $h^0\ge r+1$. \Cref{tab:A} lists, for every level, the argument for each of the two assertions of
Theorem~A.

\subsection{The framework}
\begin{lemma}\label{lem:inject}
If $C$ has no function of degree $5$ over $\Q$, then the points of degree $5$ inject into $\W(\Q)$, and a
point of degree $5$ is the unique effective divisor in its class.
\end{lemma}
\begin{proof}
If $D$ is a point of degree $5$ and $h^0(D)\ge2$, a nonconstant $f\in L(D)$ over $\Q$ has a nonzero rational
pole divisor bounded by $D$; since $D$ is irreducible, it equals $D$, so $\deg f=5$.
\end{proof}

Up to isogeny, $J$ is a product of modular abelian varieties. Let $\A\subseteq J$ be the sum of the isotypic
components whose $L$-function vanishes at $s=1$; the vanishing is decided by exact modular-symbol
computations. Each isotypic factor of $J/\A$ has $L(1)\ne0$, so $(J/\A)(\Q)$ is finite by Kato
\cite[Corollary~14.3]{Kato2004} (see also Kolyvagin and Logachev \cite{KolyvaginLogachev1989}). By Faltings
\cite{Faltings1991}, the Zariski closure of $\W(\Q)$ is a finite union of translates of abelian subvarieties;
a positive-dimensional component $T=x+B'$ with dense rational points is defined over $\Q$, and its direction
$B'$ lies in $\A$, because the image of $T$ in $\Pic^5(C)/\A$ is a translate of an abelian variety with dense
rational points in a variety with finitely many. So it suffices to show, in addition to excluding functions of
degree $5$, that every such $T$ has a fixed proper subdivisor on its locus with $h^0=1$: by \cref{lem:inject}
the classes of points of degree $5$ then lie in the finitely many isolated points of the closure of $\W(\Q)$.

At each working prime $p\nmid N$ Magma's canonical model has good reduction (\S\ref{sec:prelim}), and we
enumerate all effective divisors of degree $5$ over $\Fp$, for every partition of $5$ into degrees of places.
If $x+B\subseteq\W$ over $\Q$, then the reduction contains the entire coset $\bar x+B(\Fp)$ in $\W(\Fp)$.

\begin{lemma}[Local rigidity]\label{lem:local}
Let $R=\Z_p$, let $Z$ be a closed subscheme of a scheme smooth over $R$, and let $U\to Z$ be a morphism from a
smooth $R$-scheme of relative dimension $n$. Let $u\mapsto z$ be rational points of the special fibres. If
$\dim T_zZ_{\Fp}=n$ and $T_uU_{\Fp}\to T_zZ_{\Fp}$ is an isomorphism, then $U\to Z$ induces an isomorphism of
completed local rings at $u$ and $z$.
\end{lemma}
\begin{proof}
Lifting cotangent generators gives a surjection $R[[x_1,\dots,x_n]]\to\widehat{\mathcal O}_{Z,z}$; its
composite with $\widehat{\mathcal O}_{Z,z}\to\widehat{\mathcal O}_{U,u}\cong R[[t_1,\dots,t_n]]$ has invertible
linear part modulo $p$, so it is an isomorphism, and the surjection is injective.
\end{proof}
Throughout, $W^r_5$ carries its determinantal scheme structure, defined by the same minors on the relative
Picard scheme of the smooth model over $\Z_p$; so it is a closed subscheme of a smooth $\Z_p$-scheme, its special
fibre is $W^r_5$ of $C_{\Fp}$, and tangent spaces may be computed over $\Fp$. At a class $[D]$ of $\W$ with
$h^0(D)=1$ the annihilator of $T_{[D]}\W$ in $H^0(K_C)$ is $H^0(K_C-D)$, and at a class with $h^0\ge2$ the
tangent space of $\W$ is the whole tangent space of $\Pic^5$. At a class $L$ of $W^1_5$ with $h^0(L)=2$ the
tangent space of $W^1_5$ is the annihilator of the image of the Petri map
$H^0(L)\otimes H^0(K_C-L)\to H^0(K_C)$; if $L=A+P$ with $A$ base-point free of degree $4$, the base-point-free
pencil trick gives $\dim T_LW^1_5=8-g+h^0(K_C-L-A)$ \cite[Chapter~IV, \S4]{ACGH1985}.

We apply \cref{lem:local} to closed subschemes of the form $Z_S=\bigcap_{s\in S}(\W-s)$ of $\Pic^5$ over $\Z_p$,
where $S\subseteq J(\Z_p)$ is a finite set of lifts of points of $J(\Fp)$. If $n$ is prime to $p$, then $J[n]$ is
finite étale over $\Z_p$, so every point of order dividing $n$ in $J(\Fp)$ has a unique lift to $J[n](\Z_p)$, and
the same holds in every abelian subvariety of $J$ (these have good reduction at $p$, since their Tate modules are
submodules of that of $J$). In particular, if a subgroup of $J(\Fp)$ of order prime to $p$ lies in the
reductions of two abelian subvarieties, its lift lies in both.

\subsection{Excluding functions of degree 5}
For $21$ of the $30$ levels, the absence of a function of degree $5$ is already in \cite[\S7.2 and
Table~2]{DHJO}: at $83,89,92,101,107,111,131,141,142,143,167,191$ by \cite[Proposition~5.2]{DHJO}, and at
$82,86,99,118,123,128,145,155,159$ because the $\Q$-gonality is at least $6$ \cite{NajmanOrlic2024}. It is new
at $74,77,79,85,88,103$ and at $96,98,100$. For completeness we give the argument we use at every level.

\emph{The double-cover inequality.} If $C$ has an involution $w$ with quotient of genus $h$ and $g>2h+4$, a
morphism of degree $5$ cannot exist, even over $\overline\Q$, by \cref{lem:CS} (the degrees $2$ and $5$ are
coprime). This applies to $17$ levels (\cref{tab:A}). At $142,143,167,191$, \cite[Proposition~5.2]{DHJO} argue
instead with a morphism of degree $4$ to $\PP^1$: \cref{lem:CS} with degrees $4$ and $5$ gives $g\le12$, while
these curves have genus at least $13$.

\emph{Published gonality.} $\gon_\Q X_0(82)=6$ \cite[Theorem~1.4]{NajmanOrlic2024}.

\emph{Empty special loci.} At $N=86,99,128$ and $p=7,5,3$ respectively, the numbers of effective divisors of
degree $5$ over $\Fp$ and of their classes are equal, so $W^1_5(\Fp)=\varnothing$, and \cref{lem:specialization}
applies.

\emph{Local pencils.} At $N=77,85,88$ the curve has $\Q$-gonality $4$ \cite[Theorem~1.2]{NajmanOrlic2024}; let
$A$ be the class of a rational pencil of degree $4$. At $p=5,7,3$ respectively, the certificate checks that every
effective divisor $D$ of degree $5$ over $\Fp$ with $h^0(D)\ge2$ has $h^0(D)=2$, a base point, and moving part of
degree $4$; that these moving parts form a single class $\bar A_4$; and that $\dim T_{[D]}W^1_5=8-g+h^0(K-D-\bar A_4)=1$.
The reduction $\bar A$ of $A$ has degree $4$ and $h^0\ge2$ (\cref{lem:specialization}). For a rational point
$\bar P$, the class $\bar A+\bar P$ has $h^0\ge2$, so by the checks $h^0(\bar A+\bar P)=2$ and its moving part has
degree $4$; hence $h^0(\bar A)=2$ and $\bar A$ is base-point free (a base point of $\bar A$ would make the moving part
of $\bar A+\bar P$ of degree at most $3$). So $\bar A$ is the moving class $\bar A_4$, and $h^0(A)=2$.
Let $L$ be a rational class in $W^1_5$. Its reduction lies in $W^1_5(\Fp)$, so it is $\bar A_4+\bar P$ for a point
$\bar P\in C(\Fp)$. Apply \cref{lem:local} to $U=C\to Z=W^1_5$, $P\mapsto A+P$, at $\bar P\mapsto\bar A_4+\bar P$: the
Abel map is an immersion, and both tangent spaces have dimension $1$. So near $\bar A_4+\bar P$ every $\Z_p$-point
of $W^1_5$ is $A+P$ with $P\in C(\Z_p)$; in particular $L=A+P$, and $P=L-A$ is a rational effective class of
degree $1$, that is, a rational point. Since $h^0(L)=2=h^0(A)$, all sections of $L$ vanish at $P$. So every rational
class in $W^1_5$ has a base point, and there is no function of degree $5$.

Level $74$ uses the same argument at the two primes $3$ and $5$, after a sieve. Here $\gon_\Q X_0(74)=4$
\cite[Theorem~1.2]{NajmanOrlic2024}; a rational pencil of degree $4$ is the pullback of the hyperelliptic pencil of
$X_0(74)/w_{74}$, which has genus $2$. The positive-rank part $\A$ is isogenous to $E_{37}^2$, where $E_{37}$ is
the elliptic curve $37a1$. The rank-zero quotient $R=J/\A$ has $R(\Q)$ finite, and $\#R(\Q)$ divides $\#R(\F_q)$ for
every odd prime $q\ne37$ \cite[Appendix]{Katz1981}; the gcd over the odd primes $q\le31$ is $513$. Since
$\#E_{37}(\F_3)=7$ and $\#E_{37}(\F_5)=8$, $\A(\F_3)$ is a $7$-group and $\A(\F_5)$ a $2$-group, so
$\A(\Q)_{\mathrm{tors}}$, which injects into both, is trivial. The differences of the four rational cusps generate a
finite subgroup $\mathcal C\subseteq J(\Q)$ (Manin--Drinfeld \cite{Manin1972,Drinfeld1973}) whose reductions modulo
$3$ and $5$ have order $513$. It meets the torsion-free group $\A(\Q)$ trivially, so it maps injectively, hence
isomorphically, onto $R(\Q)$, and $J(\Q)=\mathcal C+\A(\Q)$. Every rational class of degree $5$ is therefore
$[5\infty]+z+a$ with $z\in\mathcal C$ and $a\in\A(\Q)$, and its reduction lies in $[5\infty]+\bar z+J(\Fp)[\ell^\infty]$,
where $\ell=7$ for $p=3$ and $\ell=2$ for $p=5$ (these primary subgroups have orders $49$ and $256$, and contain
$\A(\Fp)$). Call $z\in\mathcal C$ \emph{exceptional at $p$} if some class in $[5\infty]+\bar z+J(\Fp)[\ell^\infty]$
lies in $W^1_5(\Fp)$ and either has a base-point-free pencil of degree $5$ or has $\dim T W^1_5\ne1$. The certificate
(\texttt{pencils74.m}; see also \cite{SuppA}) finds exactly two exceptional $z$ at each prime, and the two pairs are disjoint. For every $z$ we choose a prime at
which $z$ is not exceptional; there the argument of the previous paragraph applies verbatim to every rational class
$[5\infty]+z+a$ in $W^1_5$ (the certificate also checks that $W^1_5(\Fp)$ has a unique moving class of degree $4$).
So $X_0(74)$ has no function of degree $5$.

The remaining levels are treated with the finiteness argument below: rank-zero lifting at $96,98,100$ and the
quotient sieve at $79,103$.

\subsection{Route R: direct reduction}
\emph{Rank zero ($96,98,100$).} Here $L(J,1)\ne0$, so $J(\Q)$ is finite \cite[Corollary~14.3]{Kato2004} and
reduction is injective on it at the odd good primes $p=7,3,3$ \cite[Appendix]{Katz1981}. This is the situation of
\cite[Proposition~7]{DerickxVanHoeij2014}. The computation
shows that every special class of degree $5$ over $\Fp$ has $h^0=2$ and is $\bar A_4+\bar P$ for a class $\bar
A_4$ of degree $4$ with two sections and a point $\bar P\in C(\Fp)$, and it lifts every such $\bar A_4$ to a
rational class $A_4$. At $96$ and $100$ every point of $C(\Fp)$ is the reduction of a rational point $P$, and
injectivity forces a rational special class to be $A_4+P$, which has the base point $P$. At $98$ only four of the
six points of $C(\F_3)$ (the cusps) lift. If $L$ is a rational special class, then $L-A_4$ is a rational class of
degree $1$ whose reduction is $\bar P$; its push-forward along the degeneracy map $X_0(98)\to X_0(49)=E$ is a point
of $E(\Q)\cong\Z/2$ reducing to the image of $\bar P$. For a noncuspidal $\bar P$ that image is a noncuspidal
point of $E(\F_3)\cong\Z/4$ (degeneracy maps send noncusps to noncusps), of order $4$, which is not the
reduction of a point of $E(\Q)$. So $\bar P$ is a cusp, and $L=A_4+P$ with a fixed rational cusp $P$.

\emph{Empty cosets ($85,88,92$).} Here $\A$ is simple, of dimension $2,1,1$. At $p=3$, the $\ell$-primary
subgroup of $\A(\F_3)$, for $\ell=17,7,7$, has no full coset inside $\W(\F_3)$; so no translate exists.

\subsection{Route Q: a rank-zero quotient (nine prime levels)}
For $N\in\{79,83,89,\allowbreak 101,103,107,\allowbreak 131,167,191\}$ let $Y=C/w_N$, $A=\pi^*\Jac Y$ and $R=J/A$; here $\A=A$ is simple
of multiplicity one and $R$ has rank $0$. Let $t=[w_N(\infty)-\infty]\in J(\Q)$. The gcd $m$ of $\#R(\F_q)$ over
the odd primes $q\le19$ bounds $\#R(\Q)$, and at the working prime the image of $t$ in $J(\Fp)/A(\Fp)$ has order $m$
(the certificate checks $\#\langle\bar t,A(\Fp)\rangle=m\cdot\#A(\Fp)$). So $R(\Q)$ is cyclic of order $m$, generated by
the image of $t$, and reduction distinguishes the $m$ cosets $it+A(\Q)$, $0\le i<m$. Every rational class of degree $5$ lies on some $T_i=[5\infty]+it+A$, $0\le i<m$, and its reduction lies in
$[5\infty]+i\bar t+A(\Fp)$. If $T_i\subseteq\W$, then the whole coset $[5\infty]+i\bar t+A(\Fp)$ lies in $\W(\Fp)$,
and only the following indices pass this test: $i\in\{1,2,3,4\}$ for $N=79,89,101,131$; $\{1,2,3,4,23\}$ for $N=83$;
$\{2,3\}$ for $N=103,107,167,191$. For $g(Y)=1$ these are the families
$\pi^*\Pic^1(Y)+3P$, $\pi^*\Pic^2(Y)+P$, $\pi^*\Pic^2(Y)+Q$, $\pi^*\Pic^1(Y)+3Q$ with $P=\infty$,
$Q=w_N(\infty)$, and for $g(Y)=2$ only the middle two. The cusp is a fixed subdivisor: if $D$ contains no full
fibre of $\pi$ and $d+\deg D<g-2g(Y)+1$, then $H^0(C,\pi^*M(D))=H^0(Y,M)$ for every $M$ of degree $d$. Indeed,
$\mathcal O_Y\subseteq\pi_*\mathcal O_C(D)$ is saturated because $D$ contains no full fibre, and the quotient is a line
bundle of degree $\deg D-(g-2g(Y)+1)$; after tensoring with $M$ it has no sections. At
$N=83$ the index $23$ is the family $\pi^*\Pic^1(Y)+D$ through an explicit $w_{83}$-fixed cubic point $D$: with
$F=\infty+w_{83}(\infty)$ and $c=[F+D-5\infty]$, invariance of $F+D$ under $w_{83}$, which acts as $-1$ on $R$,
gives $2c=5t$ in $R(\Q)\cong\Z/41$, so $c=23t$; the fixed subdivisor is $D$, which contains no full fibre of $\pi$:
its three points are fixed by $w_{83}$ (checked in \texttt{cubic83.m}), so they are ramification points, whose fibres
are $2Q$, and $D$ is reduced. At $79$ and $103$ the certificate also checks that every class of $W^1_5(\Fp)$ in
$[5\infty]+\langle\bar t\rangle+A(\Fp)$ lies over the index $2$ or $3$. A rational class $L$ in $W^1_5$ reduces into
this set, so it lies on $T_2$ or $T_3$; these are the families $\pi^*\Pic^2(Y)+P$ and $\pi^*\Pic^2(Y)+Q$, on which
all sections vanish at the cusp $P$ or $Q$. So $L$ has a base point, and there is no function of degree $5$.

\subsection{Route T: known families (fifteen levels)}
For the remaining levels, $C$ has a double cover $\pi\colon C\to Y$ with $g(Y)=2$, or a map $\pi$ of degree $4$ to
a curve $Y$ of genus $1$ (\cref{tab:routeT}). Let $B=\pi^*\Jac Y\subseteq J$ and let $F_4$ be a fibre of $\pi$ in the
second case and twice a fibre in the first. Every class in $B+[F_4]$ is $\pi^*$ of a class of degree $2$ or $1$ on
$Y$, hence effective, so
\[
 V=B+[F_4]+C=\{\,b+[F_4]+P: b\in B,\ P\in C\,\}\subseteq\W ,
\]
the image of the proper map $B\times C\to\Pic^5(C)$; it is closed, of dimension $\dim B+1$, and stable under
translation by $B$.

\emph{The certificate} (see also \cite{SuppA}). At the prime $p$ of \cref{tab:routeT}, a set of subgroups $S\subseteq J(\Fp)$ is tested: for
a list of primes $\ell\ne p$ (and $n=8$ at $N=86$), all subgroups of $J(\Fp)$ of order $\ell$, chosen so that every
simple factor $A_i$ of $\A$ has $\#A_i(\Fp)$ divisible by some tested $\ell$. For every tested $S$ and every full
coset $x+S\subseteq\W(\Fp)$ the certificate checks
\begin{enumerate}
\item[(i)] $S\subseteq B(\Fp)$;
\item[(ii)] $x+S\subseteq B(\Fp)+[F_4]+C(\Fp)$;
\item[(iii)] $\dim T_xZ_S=\dim B+1$, where $T_xZ_S$ is computed as the common annihilator of the spaces
 $H^0(K_C-D)$ for the classes $x+s=[D]$, $s\in S$, with $h^0(D)=1$;
\item[(iv)] the differential of $B\times C\to\Pic^5(C)$ is injective at every point of $B(\Fp)\times C(\Fp)$.
\end{enumerate}
At $82$, $99$ and $128$ the positive-rank part $\A$ is a single elliptic curve of multiplicity one, which is $B$
(identified through its newform); there the only test is $S=B(\Fp)$.

\emph{Every translate lies in $V$.} Let $T=t+B'\subseteq\W$ be a positive-dimensional translate with $T(\Q)$ dense,
as in the framework, with $t\in T(\Q)$; then $B'\subseteq\A$. Suppose first that we are not at $82,99,128$.
The abelian variety $B'$ contains a simple abelian subvariety isogenous to some $A_i$, so $\#B'(\Fp)$ is divisible by
a tested prime $\ell$, and $B'(\Fp)$ contains a subgroup $S$ of order $\ell$, which is tested. (At $86$, $\A$ is
isogenous to $E_{43}^2$ with $E_{43}=43a1$, the only curve in its isogeny class, so $B'$ contains a copy of $E_{43}$,
and $E_{43}(\F_7)\cong\Z/8$ gives a cyclic subgroup of order $8$.) The reduction of $T$ is the translate
$\bar t+B'_{\Fp}$, which lies in $(\W)_{\Fp}$, so $\bar t+S$ is a full coset in $\W(\Fp)$ and (i)--(iv) hold for it.
Let $\tilde S\subseteq J(\Z_p)$ be the lift of $S$; as $\ell\ne p$, it lies in both $B'$ and $B$ (by (i)). Hence
$T+s=T$ and $V+s=V$ for $s\in\tilde S$, so $T$ and $V$ both lie in $Z_{\tilde S}=\bigcap_{s\in\tilde S}(\W-s)$. At
$82,99,128$ we have $B'=B$; we lift the points of $S=B(\Fp)$ to arbitrary points of $B(\Z_p)$ (possible since $B$ is
smooth), and again $T,V\subseteq Z_{\tilde S}$, because both are stable under all of $B$.

By (ii), $\bar t=\bar b+[F_4]+\bar P$ with $\bar b\in B(\Fp)$ and $\bar P\in C(\Fp)$; lift $(\bar b,\bar P)$ to a point
$u$ of $B\times C$ over $\Z_p$. Apply \cref{lem:local} to $B\times C\to Z_{\tilde S}$ at $\bar u\mapsto\bar t$: by
(iii) and (iv) the tangent map is an isomorphism onto a space of dimension $\dim B+1$. So the completed local ring of
$Z_{\tilde S}$ at $\bar t$ is that of $B\times C$, and every $\Z_p$-point of $Z_{\tilde S}$ reducing to $\bar t$ lies in
$V$. The points of $T(\Q_p)=T(\Z_p)$ reducing to $\bar t$ form a nonempty $p$-adic open subset of the smooth
irreducible variety $T$, which is Zariski dense in $T$; as $V$ is closed, $T\subseteq V$.

\emph{No translate contributes points of degree $5$.} Let $q\colon\Pic^5(C)\to\Pic^5(C)/B$ be the quotient, a torsor
under $J/B$. Then $q(V)$ is the image of the curve $C$, whose differences generate $J/B$, and $q(T)$ is a translate of
the abelian variety $q(B')$. In every row of \cref{tab:routeT}, $\dim(J/B)=g-\dim B>1$. If $q(B')\ne0$, then $q(T)$
would be a curve, hence equal to the irreducible curve $q(V)$, and the differences of points of $q(V)$ would generate
only the one-dimensional $q(B')$; so $q(T)$ is a point. Hence $T\subseteq B+[F_4]+P$ for a single point
$P\in C(\overline\Q)$ ($T$ is irreducible and the fibres of $C\to q(V)$ are finite). Since $T$ is defined over $\Q$,
also $T\subseteq B+[F_4]+\sigma P$ for every $\sigma\in\operatorname{Gal}(\overline\Q/\Q)$.
By the exclusion of functions of degree $5$ and \cref{lem:inject}, a class of a point of degree $5$ has $h^0=1$.
If $T\subseteq W^1_5$, then $T$ contains no such class. Otherwise $T\setminus W^1_5$ is a nonempty open subset of
$T$, so its rational points are dense in $T$. Let $L\in T(\Q)$ with $h^0(L)=1$, and let $D_L$ be its unique effective
divisor. As $L$ has an effective
representative of the form $\pi^*(e)+\sigma P$ for each $\sigma$, $D_L$ contains the Galois orbit $O$ of $P$, so
$\deg O\le5$. If $\deg O=5$, then $D_L=O$ for all such $L$; these $L$ are dense in $T$, so $T$ would be a point. Hence
$O$ is a rational effective divisor of degree at most $4$ contained in every $D_L$, and no $D_L$ is a point of
degree $5$.

\begin{proof}[Proof of Theorem A]
Combine the function exclusions with the three routes, as in \cref{tab:A}. The $36$ certificates are in the
accompanying repository (\texttt{certificates/A\_x0\_quintic}), with a list of the data per level; the
supplementary write-up \cite{SuppA} gives the computation at each level in more detail.
\end{proof}

\begin{table}[ht]\centering\small
\caption{Theorem A: the thirty levels. The route proves finiteness of the points of degree $5$ assuming there is
no function of degree $5$; the last column is the separate argument excluding functions of degree $5$ (``double
cover'': the inequality $g>2h+4$).}\label{tab:A}
\begin{tabular}{rrll@{\qquad}rrll}\toprule
$N$&$g$&route&degree-5 functions&$N$&$g$&route&degree-5 functions\\\midrule
74&8&T&local pencil, 2 primes&107&9&Q&double cover\\
77&7&T&local pencil&111&11&T&double cover\\
79&6&Q&quotient sieve&118&14&T&double cover\\
82&9&T&\cite{NajmanOrlic2024}&123&13&T&double cover\\
83&7&Q&double cover&128&9&T&$W^1_5(\F_3)=\varnothing$\\
85&7&R&local pencil&131&11&Q&double cover\\
86&10&T&$W^1_5(\F_7)=\varnothing$&141&15&T&double cover\\
88&9&R&local pencil&142&17&T&double cover\\
89&7&Q&double cover&143&13&T&double cover\\
92&10&R&double cover&145&13&T&double cover\\
96&9&R&rank zero&155&15&T&double cover\\
98&7&R&rank zero&159&17&T&double cover\\
99&9&T&$W^1_5(\F_5)=\varnothing$&167&14&Q&double cover\\
100&7&R&rank zero&191&16&Q&double cover\\
101&8&Q&double cover&&&&\\
103&8&Q&quotient sieve&&&&\\\bottomrule
\end{tabular}
\end{table}
